\chapter{ЗАДАЧА ТЕРМОПРУЖНОСТІ У М'ЯКИХ ТКАНИНАХ ТІЛА ЛЮДИНИ}

%Hidden cite for correct bibliography ordering
\nocite{bahvalov-et-al,benerdge-et-al} 

\section{Огляд проблеми}

Праця \cite{lung-tumor-thermal-conductivity} ??? наводить приклад інструменту, який може передавати температуру від 
м'якої тканини пацієнта до хірурга. Такий інструмент усуває ще один недолік пов'язаний із малоінвазивною хірургією. 
Відсутність тактильного відчуття температури можна компенсувати наявністю теплового ендоскопа, як будо описано в одному 
із попередніх розділів. Проте такий тип інформації перевантажує візуалізацію і може вимагати ще одного інструмента у 
тілі пацієнта. Що у свлю чергу призводить до потреби у додатковому порті (отворі) в оперованій порожнині. Це зумовлює
кілька наслідків пов'язаних із цим:
\begin{itemize}
    \item Більше післяопераціних болів.
    \item Довший час одужування.
    \item Більша ймовірність пошкодити судини під час пробивання додаткового порту.
\end{itemize}

\noindent Рішенням для вищезгаданих проблем міг би стати керуючий пристрій із тактильно температурним відшуклом для 
хірурга. Такий інструмент наблизив би можливості відчуттів хірурга до відкритих операцій.

--------

Одним із нововведень у малоінвазивній робототехніці був тактильний відгук від інструментів. Це була суттєва різниця між
відкритою операцією і малоінвазивною хірургією. Під час вікритої операції хірург міг\textbackslashмогла відчути певні
характеристики руками і оцінити властивосте м'яких тканин. Тільки остання версія робота хірурга від Intuitive Machines
да Вінчі отримав таку функціональність (???джерело???). Першим роботом був Senhance із такою функціональністю
(джурело??? Стання на NHI). Тим не менше, наразі обмежене коло роботів хірургів надає таку функціональність.

Для того щоб можна було розробляти такі системи управління роботами із тактильним відгуком потрібно створити симульоване 
середовище, яке буде надавати достовірні дані для калібрування інструментів. Його також можна використати для навчання
хірургів і виробляння звички ефективного управління такою системою. Однією із переваг такого керування роботом була б 
здатність відчувати неоднорідні включення у м'яких тканинах, які по аналогії із відкритою операцією можна використати 
для виявлення пухлин чи інших злоякісних утворень.

У цьому розділі розглянуто дві задачі. Перша плоска, а друга - просторова стаціонарна задача контакту кінця хірургічного 
інструменту із неоднорідним включенням у м'яких тканинах. Подається візуалізація відповідних сил.

\section{Плоска стаціонарна задача}

\subsection{Формулювання задачі}

Розглянемо область у якій є дві частини (рис. \ref{fig:elasticity_2d_domain}). $\Omega_1$ відповідає м'яким тканинам, а
$\Omega_3$ - це злоякісне утворення. Область $\Omega_2$ - це хіргучний інструмент, який контактує із м'якими тканинами
по межі $\Gamma_2$, тобто $\Gamma_2=\partial\Omega_2 \cap \partial\Omega_1$, 
$\Gamma_1 = \partial \Omega_1 \backslash \Gamma_2$.

\begin{figure}[ht!]
    \centering
    \begin{tikzpicture}
        \draw (5.7,3) rectangle (6.3,7);
        \draw (1,1) rectangle (11,3);
        \draw (6,2.5) circle (0.5);
        \draw[thick, ->] (5.3,6) -- (5.3,4);

        \draw[thick,->] (0,0) -- (12,0) node[anchor=north west] {x};
        \draw[thick,->] (0,0) -- (0,9) node[anchor=south east] {y};        
    
        \foreach \x in {0,...,4}
            \pgfmathtruncatemacro{\label}{\x * 10}
            \draw ( \x * 2.5 cm,1pt) -- ( \x * 2.5 cm,-1pt) node[anchor=north] {\label};
        \foreach \y in {0,...,3}
            \pgfmathtruncatemacro{\label}{\y * 10}
            \draw (1pt,\y * 2.5 cm) -- (-1pt,\y * 2.5 cm) node[anchor=east] {\label};
        
        \node at (11.5,2) {a};
        \node at (6,0.5) {b};
        \node at (6,7.5) {c};
        \node at (6,3.5) {$\Gamma_2$};

        \node at (8,2) {$\Omega_1$};
        \node at (6,5) {$\Omega_2$};
        \node at (6,2.5) {$\Omega_3$};
        \node at (4,0.5) {$\Gamma_1$};

    \end{tikzpicture}
    
    \caption{Переріз тканин і інструмента}
    \label{fig:elasticity_2d_domain}
\end{figure}

\noindent Параметр a рівний 8 см, b --- 40 см, а товщина хірургічного інструменту c --- 6 мм. Радіус кола $\Omega_3$ r 
дорівнює 5 мм.

Розглянемо тіло, що займає область $\Omega$ із краєм $\Gamma$. Нехай коефіцієнт теплопровідності $\lambda_{t}(x)$ та
коефіцієнти Ляме $\lambda(x),\mspace{9mu}\mu(x)$ матеріалу тіла описуються неперервними функціями, які набувають сталих
значень скрізь у $\Omega$ за винятком локальних областей неоднорідності $\Omega_{m}\subset\Omega$ ($m =
\overline{1,M}$), причому ${\Omega_{m} \cap \Omega_{l}} = \varnothing$($m \neq l$), де вони залежать від декартових
координат ${x = (}x_{1},x_{2})$, а саме 
$\lambda_{t}(x{) = {\lambda_{t0} + {\sum\limits_{m = 1}^{M}{\lambda_{\mathit{\text{tm}}}(x{) \cdot \chi_{m}}(x)}}}}$,
$\mu(x{) = {\mu_{0} + {\sum\limits_{m = 1}^{M}{\mu_{m}(x{) \cdot \chi_{m}}(x)}}}}$,
$\lambda(x{) = {\lambda_{0} + {\sum\limits_{m = 1}^{M}{\lambda_{m}(x{) \cdot \chi_{m}}(x)}}}}$.(5.1)

Тут $\chi_{m}(x)$ -- характеристична функція області $\Omega_m$ $\lambda_{t0}$, -- коефіцієнти в $\Omega{}_{m =
1}^{M}\Omega_{m}$; ${\lambda_{t0} + \lambda_{\mathit{\text{tm}}}}(x)$, ${\lambda_{0} + \lambda_{m}}(x)$, ${\mu_{t0} +
\mu_{\mathit{\text{tm}}}}(x)$ -- коефіцієнти області Ω\textsubscript{\emph{m}}, причому
$\lambda_{\mathit{\text{tm}}}(x){|_{\Gamma_{m}} = 0}$, $\lambda_{m}(x){|_{\Gamma_{m}} = 0}$, $\mu_{m}(x){|_{\Gamma_{m}}
= 0}$, де Г\textsubscript{\emph{m}} -- край області Ω\textsubscript{\emph{m}}.

Підставивши у рівняння теплопровідності та рівняння рівноваги в переміщеннях вирази (5.1), співвідношення
Дюгамеля-Неймана

\begin{equation*}
    \begin{split}
        \sigma_{\mathit{\text{ij}}}(x{) = \lambda_{0}}\delta_{\mathit{\text{ij}}}\varepsilon_{\mathit{\text{kk}}}(x{) + 2}\mu_{0}
            \varepsilon_{\mathit{\text{ij}}}(x{) - \beta_{0}}\theta + \\ + {{\sum\limits_{m = 1}^{M}{\lbrack\lambda_{m}(x)\delta_{\mathit{\text{ij}}}
            \varepsilon_{\mathit{\text{kk}}}(x{) + 2}\mu_{m}(x)\varepsilon_{\mathit{\text{ij}}}(x)-\beta_{m}(x)\theta(x)\rbrack}} \cdot \chi_{m}}(x)
    \end{split}
\end{equation*}

\noindent та співвідношеннями Коші, одержимо систему диференціальних рівнянь:

$\mathit{\Delta\theta} = {- f_{\theta}}$,(5.2)

$\mu_{0}{u_{i,\mathit{\text{jj}}} + (}{\lambda_{0} + \mu_{0}}){u_{j,\mathit{\text{ji}}} = {{- f_{\mathit{\text{ij}},j}} + \beta\theta_{,i} + \beta_{,i}}}\theta$,(5.3)

\noindent де $\Delta$ -- оператор Лапласа,
${f_{\theta} = f_{\theta}}(\theta_{,i}{) = {\sum\limits_{m = 1}^{M}{\chi_{m}\lambda_{\mathit{\text{tm}},i}{\theta_{,i}/\lambda_{\mathit{\text{tm}}}}}}}$,
\begin{equation*}
    {f_{\mathit{\text{ij}}} = f_{\mathit{\text{ij}}}}(\varepsilon{) = {\sum\limits_{m = 1}^{M}{\lbrack\lambda_{m}\delta_{\mathit{\text{ij}}}{\varepsilon_{\mathit{\text{kk}}} + 2}\mu_{m}\varepsilon_{\mathit{\text{ij}}}\rbrack\chi_{m}}}}    
\end{equation*}
${\beta_{0} = (}2{\mu_{0} + 3}\lambda_{0}{) \cdot \alpha}$, $\beta_{m}(x{) = (}2\mu_{m}(x{) + 3}\lambda_{m}(x){) \cdot
\alpha}$; $\beta(x{) = {\beta_{0} + {\sum\limits_{m = 1}^{M}{\beta_{m}(x{) \cdot \chi_{m}}(x)}}}}$, $\sigma_{ij}$,
$\varepsilon_{ij}$ -- компоненти відповідно тензора напружень та тензора деформацій; $u_{ij}$ -- компоненти вектора
переміщень; $\theta$ -- стаціонарне температурне поле; $\alpha$ -- коефіцієнт теплового розширення матеріалу;
$\delta_{ij}$-- символ Кронекера.

\noindent Вважатимемо, що на P\textsubscript{1}, P\textsubscript{2 } задані температура й тепловий потік

$\theta{|_{P_{1}} = \theta_{p}}$,
$q{|_{P_{2}} = q_{p}}$(5.4)

\noindent а на P\textsubscript{1}, P\textsubscript{2 } -- переміщення та поверхневі зусилля

$u_{i}{|_{P_{1}} = u_{\mathit{\text{pi}}}}$,
$t_{i}{|_{P_{2}} = t_{\mathit{\text{pi}}}}$, (5.5)

\noindent де ${P_{1} \cup P_{2}} = \Gamma$, ${P_{1} \cap P_{2}} = \varnothing$, $t_i=\sigma_{ij}n_j$, $n_j$ --
компоненти одиничного вектора зовнішньої нормалі до краю $\Gamma$. Отже, для визначення температурних переміщень $u_j$
маємо незв'язану змішану крайову задачу (5.2)-(5.5).

Апроксимувавши компоненти невідомих похідної від температури $\theta_{,i}$ та тензора деформацій $\varepsilon_{ij}$ в
областях Ω\textsubscript{m} функціями $\theta_{,i}^{m}$, $\varepsilon_{\mathit{\text{ij}}}^{m}$ і використавши їх у
$f_{\theta}$, перейдемо від крайової задачі відносно $\theta$, $u_i$ до крайової задачі відносно $\theta^*$, $u^*$

$\mathit{\Delta\theta}^{} = {- f_{\theta}}$,(5.6)

$\theta^{}{|_{P_{1}} = \theta_{p}}$,\emph{}$q^{}{|_{P_{2}} = q_{p}}$,(5.7)

$\mu_{0}{u_{i,\mathit{\text{jj}}}^{} + (}{\lambda_{0} + \mu_{0}}){u_{j,\mathit{\text{ji}}}^{} = {- f_{\mathit{\text{ij}},j}}}(\varepsilon^{}) + \beta\theta_{\mathrm{,i}}^{\mathrm{}} + \beta_{,i}\theta^{}$,(5.8)

$u_{i}^{}{|_{P_{1}} = u_{\mathit{\text{pi}}}}$,
$t_{i}^{}{|_{P_{2}} = t_{\mathit{\text{pi}}}}$.(5.9)

\subsection{Дискретизація рівнянь}

Для побудови розв'язків задачі (5.6)-(5.9) застосуємо МГЕ, а також МПГЕ у поєднанні з ермітовими скінченними елементами
в областях локальних неоднорідностей. З цією метою розглянемо таку область $B\subset R^{2}$, що $\Omega\subset
B,\quad\partial{\Omega \cap \partial}{B = \varnothing}$. Позначимо приграничну область $G = B$\textsubscript{.
}Розв'язки рівнянь (5.6), (5.8) зобразимо у вигляді:

$\theta^{*\gamma}(X) = \int\limits_\gamma U(X,Y)\phi_{\theta}^{\gamma}(X)d\gamma(Y) + 
\sum\limits_{m = 1}^{M}\int\limits_{\Omega_{m}}U(X,Y)f_\theta(\theta_{,j}^{\gamma m})d\Omega_{m}(Y) + C_{\theta}$, (5.10)

$u_{i}^{\gamma}(X{) = {\int\limits_{\gamma}{E_{\mathit{\text{ij}}}(X,Y)\phi_{j}^{\gamma}(X)\mathit{d\gamma}(Y{) + {\sum\limits_{m = 1}^{M}{\int\limits_{\Omega_{m}}{E_{\mathit{\text{ij}}}(X,Y)f_{\mathit{\text{jk}},k}(\varepsilon^{\gamma_{m}}(Y))d\Omega_{m}{(Y{) -}}}}}}}}}$

${- {\int\limits_{\Omega}{E_{\mathit{\text{ij}}}(X,Y)\lbrack\beta(Y)\theta_{,j}(Y{) + \beta_{,j}}(Y)\theta(Y)\rbrack d\Omega(Y)}}} + C_{i}$,(5.11)

\noindent де $\gamma\in \{G,\Gamma\}$. Застосувавши до (5.10), (5.11) теорему Гріна інтегрування частинами,
одержимо зображення для температури та компонент вектора переміщень:

$\theta^{\gamma}(X{) = {\int\limits_{\gamma}{U(X,Y)\phi_{\theta}^{\gamma}(X)\mathit{d\gamma}(Y{) - {\sum\limits_{m = 1}^{M}{{\int\limits_{\Omega_{m}}{U_{,i}(X,Y)f_{\theta}(\theta^{\gamma_{m}}(Y))d\Omega_{m}{(Y{) + C_{\theta}}}}},}}}}}}$
(5.12)

$u_{i}^{\gamma}(X{) = {\int\limits_{\gamma}{E_{\mathit{\text{ij}}}(X,Y)\phi_{j}^{\gamma}(X)\mathit{d\gamma}(Y{) - {\sum\limits_{m = 1}^{M}{\int\limits_{\Omega_{m}}{E_{\mathit{\text{ij}},k}(X,Y)f_{\mathit{\text{jk}}}(\varepsilon^{\gamma_{m}}(Y))d\Omega_{m}{(Y{) -}}}}}}}}}$

${- {\int\limits_{\Gamma}{E_{\mathit{\text{ij}}}(X,Y)\beta(Y)\theta(Y)n_{j}(Y)\mathit{d\Gamma}(Y{) +}}}}$

$+{{\int\limits_{\Omega}{\theta(Y)\lbrack E_{\mathit{\text{ij}},j}(X,Y)\beta(Y{) - E_{\mathit{\text{ij}}}}(X,Y)\beta_{,j}(Y)\rbrack d\Omega(Y)}} + C_{i}}$(5.13)

і на основі них інтегральні зображення теплового потоку та компонент деформацій, напружень та поверхневих зусиль

$q^{\gamma}(X{) = {\int\limits_{\gamma}{Q(X,Y)\phi_{\theta}^{\gamma}(X)\mathit{d\gamma}(Y{) + {\sum\limits_{m = 1}^{M}{{\int\limits_{\Omega_{m}}{Q_{,i}(X,Y)f_{\theta}(\theta^{\gamma_{m}}(Y))d\Omega_{m}{(Y)}}},}}}}}}$
(5.14)

\begin{equation*}
    \begin{split}
        \left\{ \begin{array}{lr} \varepsilon_{ij}^{*\gamma}(X) \\ \sigma_{ij}^{*\gamma}(X) \\ t_{i}^{*\gamma} (X) \end{array} \right \} =
        \int\limits_{\gamma} \left\{ \begin{array}{lr} B_{ijk}(X,Y) \\ T_{ijk}(X,Y) \\ F_{ik}(X,Y) \end{array} \right \}
            \phi_k^\gamma(Y)d\gamma(Y) - \\ - 
        \sum \limits_{m=1}^M \int \limits_{\Omega_m} \left\{ \begin{array}{lr} B_{ijk,l}(X,Y) \\ T_{ijk,l}(X,Y) \\ F_{ik,l}(X,Y) \end{array} \right \}
            f_{kl}(\varepsilon^{\gamma m}(Y))d\Omega_m(Y) - \\ -
        \int \limits_{\Gamma} \beta(Y) \theta(Y) \left\{ \begin{array}{lr} B_{ijk}(X,Y) \\ T_{ijk}(X,Y) \\ F_{ik}(X,Y) \end{array} \right \}
            n_k(Y)d\Gamma(Y) + \\ +
        \int \limits_{\Omega} \beta(Y) \theta(Y) \left\{ \begin{array}{lr} B_{ijk,k}(X,Y) \\ T_{ijk,k}(X,Y) \\ F_{ik,k}(X,Y) \end{array} \right \}
            d\Omega(Y) - \\ - 
        \sum \limits_{m=1}^M \int \limits_{\Omega_m} \beta_{m,k}(Y) \theta(Y) 
                \left\{ \begin{array}{lr} B_{ijk}(X,Y) \\ T_{ijk}(X,Y) \\ F_{ik}(X,Y) \end{array} \right \} d\Omega_m(Y) +
        f_{ij}(\varepsilon^{\gamma m}(X)) 
        \left\{ \begin{array}{lr} 0 \\ 1 \\ n_j \end{array} \right \}, (5.15)
    \end{split}
\end{equation*}

\noindent де $U(X,Y)$ -- фундаментальний розв'язок стаціонарного рівняння теплопровідності,
$E_{\mathit{\text{ij}}}(X,Y{) = E_{\mathit{\text{ij}}} = \frac{\lambda_{0} + \mu_{0}}{4 \pi \mu_0({\lambda_{0} + \mu_{0}})}}\left\{ {{- \frac{{\lambda_{0} + 3}\mu_{0}}{\lambda_{0} + \mu_{0}}}\text{ln}{\mathit{r\delta}_{\mathit{\text{ij}}} + r_{,i}}r_{,j}} \right\}$
-- фундаментальний розв'язок системи рівнянь задачі пружності; ${Y = (}Y_{1},Y_{2}), Y_{1},Y_{2}$--- система координат,
що збігається зі системою $x_{1},x_{2}$ і використовується для опису точки, в якій діють невідомі фіктивні джерела тепла
$\phi_{\theta}^{\gamma}$ чи прикладені компоненти невідомих фіктивних поверхневих зусиль $\phi_{j}^{\Gamma}$ чи масових
сил $\phi_{j}^{G}$; $C_{\theta}$, невідомі сталі,
$B_{\mathit{\text{ijk}}}(X,Y{) = B_{\mathit{\text{ijk}}} = \frac{1}{2}}({E_{\mathit{\text{ik}},j} + E_{\mathit{\text{jk}},i}})$,
$T_{\mathit{\text{jk}}}(X,Y{) = T_{\mathit{\text{jk}}} = \lambda_{0}}\delta_{\mathit{\text{ij}}}{B_{\mathit{\text{llk}}} + 2}\mu_{0}B_{\mathit{\text{ijk}}}$,
$F_{\mathit{\text{ij}}}(X,Y{) = F_{\mathit{\text{ij}}} = T_{\mathit{\text{ijk}}}}n_{k}$.

У приграничній області \emph{G} і на ділянках краю $P_{1},P_{2}$ уведемо відповідно приграничні $G_v$ та граничні
${\Gamma_{v} = \partial}{G_{v} \cap \Gamma}\mspace{9mu}$$({v = 1},\text{.}\text{.}\text{.},V)$ елементи, при­чому
$\mathit{\text{mes}} {G_{\nu} = 2}$, ${}_{v = 1}^{V}{G_{v} = G}$, ${{G_{v} \cap G_{w}} = \varnothing},$ ${{\Gamma_{v}
\cap \Gamma_{w}} = \varnothing},$ при $v \neq w$, ${}_{v = 1}^{V_{1}}\Gamma_{v}{= P_{1}}$, ${}_{v = {V_{1} +
1}}^{V}\Gamma_{v}{= P_{2}}$. Геометрію приграничних і граничних елементів моделюємо за допомогою вектора $\Psi$ 
інтерполяційних функцій у локальній системі координат. Вздовж кожного з елементів $\Gamma_v$ та по площі елемента
$G_v$ здійснимо апроксимацію невідомих функцій $\phi_{\theta}^\gamma$, $\phi_{j}^\gamma$ за допомогою інтерполяційних 
поліномів $\phi_{\theta}^{\gamma}(\eta){|_{\gamma_{\nu}} = \psi^{T}}{(\eta{) \cdot \phi_{\theta \nu}^{\gamma}}}$,
$\phi_{j}^{\gamma}(\eta){|_{\gamma_{\nu}} = \psi^{T}}{(\eta{) \cdot \phi_{\mathit{j\nu}}^{\gamma}}}$,
де $\phi_{\theta \nu}^{\gamma}$, $\phi_{jv}^{\gamma}$ --- вектори невідомих вузлових апроксимацій функцій 
$\phi_{\theta}^{\gamma}$, $\phi_{j}^{\gamma}$ на $\nu$ -ому елементі.

Області $\Omega_m$ дискретизуємо системою ермітових скінчених елементів $\Omega_{ms}$ ($s = \overline{1,S_{m}}$, 
зобразивши функції $\theta^{\gamma m}$, $\varepsilon_{ij}^{\gamma m}$ на кожному з них пробними функціями

$\theta^{\mathit{\gamma m}}{|_{\Omega_{\mathit{\text{ms}}}} = \xi^{T}}{(\zeta_{1}}{,\zeta_{2}}{) \cdot \theta^{\gamma\mathit{\text{ms}}}}$,
$\varepsilon_{\mathit{\text{ij}}}^{\mathit{\gamma m}}{|_{\Omega_{\mathit{\text{ms}}}} = \xi^{T}}{(\zeta_{1}}{,\zeta_{2}}{) \cdot \varepsilon_{\mathit{\text{ij}}}^{\gamma\mathit{\text{ms}}}}$,

\noindent де $\theta^{\gamma_{ms}}$, $\varepsilon_{ij}^{\gamma_{ms}}$ -- вузлові вектори значень температури 
$\theta^{\gamma m}$ та компоненти $\varepsilon_{\mathit{\text{ij}}}^{\mathit{\gamma m}} $тензора деформацій на s-тому 
елементі, $\xi$ -- вектор базових функцій у локальній системі координат.

На основі таких апроксимацій одержимо дискретні аналоги формул (5.12)-(5.15)

$\theta^{\gamma}(X{) = U^{\gamma \nu}}(X{{) \cdot \phi_{\theta \nu}^{\gamma}} + U_{i}^{\mathit{\text{ms}}}}(X{{) \cdot \lambda_{t,i}^{\mathit{\text{ms}}}/\lambda_{t}^{\mathit{\text{ms}}}} + C_{\theta}}$,(5.16)

$q^{\gamma}(X{) = Q^{\gamma \nu}}(X{{) \cdot \phi_{\theta \nu}^{\gamma}} + Q_{i}^{\mathit{\text{ms}}}}(X{) \cdot \lambda_{t,i}^{\mathit{\text{ms}}}/\lambda_{t}^{\mathit{\text{ms}}}}$,(5.17)

$u_{i}^{\gamma}(X{) = E_{\mathit{\text{ij}}}^{\gamma \nu}}(X{{) \cdot \phi_{\mathit{j\nu}}^{\gamma}} + E_{i}^{\lambda\mathit{\text{ms}}}}(X{{) \cdot \varepsilon_{\mathit{\text{kk}}}^{\gamma\mathit{\text{ms}}}} + E_{\mathit{\text{ijl}}}^{\mu\mathit{\text{ms}}}}(X{{) \cdot \varepsilon_{\mathit{\text{jl}}}^{\gamma\mathit{\text{ms}}}} + C_{i}}$,
(5.18)

\begin{equation*}
    \left\{ \begin{array}{lr} \varepsilon_{ij}^{*\gamma} \\ \sigma_{ij}^{*\gamma} \\ t_{i}^{*\gamma} \end{array} \right \} (X) =
    \left\{ \begin{array}{lr} B_{ijk}^\nu \\ T_{ijk}^\nu \\ F_{ik}^\nu \end{array} \right \} (X) \phi_{k \nu}^\gamma + 
    \left\{ \begin{array}{lr} B_{ij}^{\lambda m s} \\ T_{ij}^{\lambda m s} \\ F_{i}^{\lambda m s} \end{array} \right \} (X) \varepsilon_{kk}^{\gamma m s} + 
    \left\{ \begin{array}{lr} B_{ijkl}^{\mu m s} \\ T_{ijkl}^{\mu m s} \\ F_{ikl}^{\mu m s} \end{array} \right \} (X) \varepsilon_{kl}^{\gamma m s}, (5.19)
\end{equation*}

\noindent де

$\begin{Bmatrix}
U^{\Gamma_{v}} \\
Q^{\Gamma_{v}} \\
\end{Bmatrix}{(X{) = {\int\limits_{- 1}^{1}\begin{Bmatrix}
U^{\Gamma_{v}} \\
Q^{\Gamma_{v}} \\
\end{Bmatrix}}}}\psi^{T}J_{\mathit{\Gamma\nu}}\mathit{d\eta}$,
$\begin{Bmatrix}
U^{\Gamma_{v}} \\
Q^{\Gamma_{v}} \\
\end{Bmatrix}{(X{) = {\int\limits_{- 1}^{1}\begin{Bmatrix}
U^{\Gamma_{v}} \\
Q^{\Gamma_{v}} \\
\end{Bmatrix}}}}\psi^{T}J_{\mathit{\Gamma\nu}}\mathit{d\eta}$

\begin{equation*}
    \left\{ \begin{array}{lr} E_{ij}^{\Gamma_\nu} \\ B_{ijk}^{\Gamma_\nu} \\ T_{ijk}^{\Gamma_\nu} \\ F_{ik}^{\Gamma_\nu} \end{array} \right \} (X) =
    \int \limits_{-1}^1 
        \left\{ \begin{array}{lr} E_{ij} \\ B_{ijk} \\ T_{ijk} \\ F_{ik} \end{array} \right \} \psi^T J_{\Gamma_\nu} d \eta;
    \left\{ \begin{array}{lr} E_{ij}^{G_\nu} \\ B_{ijk}^{G_\nu} \\ T_{ijk}^{G_\nu} \\ F_{ik}^{G_\nu} \end{array} \right \} (X) =
    \int \limits_{-1}^1 \int \limits_{-1}^1 
        \left\{ \begin{array}{lr} E_{ij} \\ B_{ijk} \\ T_{ijk} \\ F_{ik} \end{array} \right \} \psi^T J_{G_\nu} d \eta_1 d \eta_2;
\end{equation*}

\begin{equation*}
    \left\{ \begin{array}{lr} E_{i}^{\lambda m s} \\ B_{ij}^{\lambda m s} \\ T_{ij}^{\lambda m s} \\ F_{i}^{\lambda m s} \end{array} \right \} (X) =
    - \int \limits_{-1}^1 \int \limits_{-1}^1 
        \left\{ \begin{array}{lr} E_{ij,j} \\ B_{ijk,k} \\ T_{ijk,k} \\ F_{ik,k} \end{array} \right \} \lambda_{ms} \xi^T H_{ms} d \zeta_1 d \zeta_2 +
    \lambda_{ms} \left\{ \begin{array}{lr} 0 \\ 0 \\ \delta_{ij} \\ n_j \end{array} \right \} \xi^T \chi_{ms}(X),
\end{equation*}

\begin{equation*}
    \left\{ \begin{array}{lr} E_{il}^{\mu m s} \\ B_{ijkl}^{\mu m s} \\ T_{ijkl}^{\mu m s} \\ F_{il}^{\mu m s} \end{array} \right \} (X) =
    -2 \int \limits_{-1}^1 \int \limits_{-1}^1
        \left\{ \begin{array}{lr} E_{ij,j} \\ B_{ijk,k} \\ T_{ijk,k} \\ F_{ik,k} \end{array} \right \} \mu_{ms} \xi^T H_{ms} d \zeta_1 d \zeta_2 +
    2 \mu_{ms} \left\{ \begin{array}{lr} 0 \\ 0 \\ 1 \\ n_j \end{array} \right \} \xi^T \chi_{ms}(X),
\end{equation*}

$J_{\Gamma_{v}}$, $J_{G_{v}}$, $H_{ms}$ -- якобіани переходу від змінних $\eta$, $\eta_1$, $\eta_2$;$\zeta_1$, $\zeta_2$
до Х відповідно. $\chi_{ms}$ -- характеристична функція елемента $\Omega_{ms}$, $\lambda_{ms}$, $\mu_{ms}$ -- значення
функцій $\lambda_{m}(X)$, $\mu_{m}(X)$ при $X\in\Omega_{ms}$.

Уведемо функції нев'язок
$R_{\mathit{\text{ij}}}^{\gamma\mathit{\text{ms}}}(x{) = \varepsilon_{\mathit{\text{ij}}}^{\gamma}}(x{) - \varepsilon_{\mathit{\text{ij}}}^{\gamma\mathit{\text{ms}}}}(x)$
на $\underset{s = 1}{\overset{S_{m}}{}}\Omega_{\mathit{\text{ms}}}$ ($m = \overline{1,M}$), $R_{1i}^{\gamma}(x{) =
u_{i}^{\gamma}}(x{) - u_{\mathit{\text{pi}}}}(x)$ на $P_1$ і $R_{2i}^{\gamma}(x{) = t_{i}^{\gamma}}(x{) -
t_{\mathit{\text{pi}}}}(x)$ на $P_2$.

Унаслідок граничного переходу та дискретизації виразів для переміщень (5.9) та зусиль (5.10), одержимо їх дискретні
аналоги у точці $X_{0}{\in\underset{\nu = 1}{\overset{V}{\cup}}\Gamma_{\nu}}$:

$u_{i}^{\gamma}(X_{0}{) = E_{\mathit{\text{ij}}}^{\gamma \nu}}(X_{0}{{) \cdot \phi_{\mathit{j\nu}}^{\gamma}} + E_{i}^{\lambda\mathit{\text{ms}}}}(X_{0}{{) \cdot \varepsilon_{\mathit{\text{kk}}}^{\gamma\mathit{\text{ms}}}} + E_{\mathit{\text{ijl}}}^{\mu\mathit{\text{ms}}}}(X_{0}{{) \cdot \varepsilon_{\mathit{\text{jl}}}^{\gamma\mathit{\text{ms}}}} + C_{i}}$,
(5.11)

$t_{i}^{\Gamma}(X_{0}{) = {\pm \frac{1}{2}}}\delta_{\mathit{\text{ij}}}\psi^{T}(X_{0}{{) \cdot \phi_{\mathit{j\nu}}^{\Gamma}} + F_{\mathit{\text{ij}}}^{\mathit{\Gamma\nu}}}(X_{0}{{) \cdot \phi_{\mathit{j\nu}}^{\Gamma}} + F_{i}^{\lambda\mathit{\text{ms}}}}(X_{0}{{) \cdot \varepsilon_{\mathit{\text{kk}}}^{\Gamma\mathit{\text{ms}}}} + F_{\mathit{\text{ijl}}}^{\mu\mathit{\text{ms}}}}(X_{0}{) \cdot \varepsilon_{\mathit{\text{jl}}}^{\Gamma\mathit{\text{ms}}}}$,
$t_{i}^{G}(X_{0}{) = F_{\mathit{\text{ij}}}^{\mathit{G\nu}}}(X_{0}{{) \cdot \phi_{\mathit{j\nu}}^{G}} + F_{i}^{\lambda\mathit{\text{ms}}}}(X_{0}{{) \cdot \varepsilon_{\mathit{\text{kk}}}^{\mathit{\text{Gms}}}} + F_{\mathit{\text{ijl}}}^{\mu\mathit{\text{ms}}}}(X_{0}{) \cdot \varepsilon_{\mathit{\text{jl}}}^{\mathit{\text{Gms}}}}$(5.12)

\noindent причому точка X\textsubscript{0} така, що в ній існує єдина дотична до $\cup_{\nu = 1}^{V}\Gamma_{\nu}$;
інтеграли $F_{\mathit{\text{ij}}}^{\mathit{\Gamma\nu}}(X_{0})$ слід розуміти у сенсі головного значення Коші, а всі
решта --- у сенсі Рімана.

Для визначення невідомих векторів вузлових значень фіктивних поверхневих зусиль $\phi_{j\nu}^\Gamma$ та масових сил
$\phi_{j\nu}^G$, векторів вузлових значень компонент тензора $\varepsilon_{ij}^{\gamma ms}$ в областях неоднорідностей
та вектора $C(C_1, C2)$ побудуємо систему лінійних алгебричних рівнянь за методом зважених нев'язок. Враховуємо для
цього і дискретні аналоги

${{W^{\gamma \nu} \cdot \phi_{\mathit{j\nu}}^{\gamma}} + {W_{j}^{\lambda\mathit{\text{ms}}} \cdot \varepsilon_{\mathit{\text{kk}}}^{\gamma\mathit{\text{ms}}}} + {W_{l}^{\mu\mathit{\text{ms}}} \cdot \varepsilon_{\mathit{\text{jl}}}^{\gamma\mathit{\text{ms}}}}} = 0$
(5.13)

\noindent умов

$\int\limits_{\gamma}{\phi_{j}^{\gamma}(Y)\mathit{d\gamma}(Y{) + {\sum\limits_{m = 1}^{M}{\int\limits_{\Omega_{m}}{f_{\mathit{\text{jk}},k}(\varepsilon^{\gamma_{m}}(Y))d\Omega_{m}{(Y{) = 0}}}}}}}$,

\noindent де

$W^{\mathit{\Gamma\nu}} = {\int\limits_{- 1}^{1}{\psi^{T}J_{\mathit{\Gamma\nu}}\mathit{d\eta}}}$,
$W^{\mathit{G\nu}} = {\int\limits_{- 1}^{1}{\int\limits_{- 1}^{1}{\psi^{T}J_{\mathit{G\nu}}\mathit{d\eta}_{1}\mathit{d\eta}_{2}}}}$,
$W_{j}^{\lambda\mathit{\text{ms}}} = {\int\limits_{- 1}^{1}{\int\limits_{- 1}^{1}{\lambda_{\mathit{\text{ms}}}\xi^{T}H_{\mathit{\text{ms}}}\mathit{d\zeta}_{1}\mathit{d\zeta}_{2}}}}$,
${W_{j}^{\mu\mathit{\text{ms}}} = 2}{\int\limits_{- 1}^{1}{\int\limits_{- 1}^{1}{\mu_{\mathit{\text{ms}}}\xi^{T}H_{\mathit{\text{ms}}}\mathit{d\zeta}_{1}\mathit{d\zeta}_{2}}}}$.

У матричній формі система для знаходження невідомих $\phi_{jv}^\gamma$, $\varepsilon_{ij}^{\gamma ms}$ та $C_j$ набуде
такого загального вигляду

$\begin{bmatrix}
\begin{matrix}
E^{\gamma} \\
F^{\gamma} \\
B^{\gamma} \\
W^{\gamma} \\
\end{matrix} & {\begin{matrix}
E^{1} \\
F^{1} \\
B^{1} \\
W^{1} \\
\end{matrix}\text{.}\text{.}\text{.}\begin{matrix}
E^{M} \\
F^{M} \\
B^{M} \\
W^{M} \\
\end{matrix}} & \begin{matrix}
I \\
0 \\
0 \\
0 \\
\end{matrix} \\
\end{bmatrix}{\begin{bmatrix}
\phi^{\gamma} \\
\varepsilon^{1} \\
 \vdots \\
\varepsilon^{M} \\
C \\
\end{bmatrix} = \begin{bmatrix}
u^{p} \\
t^{p} \\
0 \\
0 \\
\end{bmatrix}}$,(5.14)

\noindent де $\phi^\gamma$ --- вектор невідомих вузлових значень фіктивних поверхневих зусиль або масових сил у всьому
дискретному аналогові $\underset{\nu = 1}{\overset{V}{}}\gamma_{v}$ краю тіла або приграничної до нього області;

$\varepsilon^m$ --- вектор невідомих вузлових значень компонент деформацій та їх похідних у дискретному аналогові
$\underset{s = 1}{\overset{S_{m}}{}}\Omega_{\mathit{\text{ms}}}$ ($m = \overline{1,M}$) області локальної неоднорідності
$\Omega_m$;

u\textsuperscript{p}, t\textsuperscript{p} --- вектори значень заданих на краю тіла переміщень та поверхневих зусиль.

Елементи блоків глобальної матриці системи є сумами внесків окремих елементів (5.10)-(5.13). А саме блок $E^\gamma$
сформований з $E_{ij}^{\gamma \nu}$; блок $E^m$ --- з $E_{i}^{\lambda ms}$ та $E_{ijl}^{\mu ms}$; блок $F^\gamma$ --- з
$F_{ik}^{\gamma \nu}$; блок $F^m$ --- з $F_{i}^{\lambda ms}$ та $F_{ijl}^{\mu ms}$; блок $B^{\gamma}$--- з
$B_{\mathit{\text{ijk}}}^{\gamma \nu}$; блок $B^m$ --- з $B_{ij}^{\lambda ms}$ та $B_{ijkl}^{\mu ms}$; блок
$W^{\gamma}$--- з $W^{\gamma \nu}$; блок $W^{m}$ --- з $W_{j}^{\lambda ms}$ та $W_{l}^{\mu ms}$.

У довільній точці $x \in \Omega \cup \Gamma$ переміщення, деформації, напруження та поверх­неві зусилля визначають з
формул (5.9), (5.10), використовуючи розв'язок системи (5.14).

\subsection{Розв'язок}

\section{Просторова стаціонарна задача}

\subsection{Формулювання задачі}

\subsection{Дискретизація рівнянь}

\subsection{Розв'язок}

\section{Висновок}

Проведені дослідження показали, що МПГЕ забезпечує вищу точність розрахунків напружено-деформованого стану у
порівнянні з МГЕ при використанні однакової кількості елементів та однакового ступеня апроксимації невідомих функцій
``фіктивних'' масових сил. Це обґрунтовується тим, що пригранична область згладжує вплив введених у ній функцій.

Точність обчислень при знаходженні напружено-деформованого стану залежить від вдалого вибору форми приграничних
елементів та значення висоти $h$ області $G^{m}$, натомість зміна шаблону інтегрування суттєвого значення не
мала. Рекомендації щодо вибору параметра $h$ можуть бути наступними: для модельних задач з відомими аналітичними
розв'язками варто виходити з апостеріорних оцінок отриманих числових результатів, для інших -- порівнювати результати
обчислень шуканих функцій у певних ділянках краю із заданими на них граничними умовами, а також точність задоволення
умов контакту на лініях контакту.

Відзначимо, що зростання (спадання) модуля зсуву в області $\Omega_{2}$ спричинює вищі (нижчі) значення компоненти
$u_{2}^{(G)}$, а зростання (спадання) коефіцієнта Пуассона -- менші (більші) значення цієї компоненти при дослідженні
кусково-однорідного тіла.

Порівняльний аналіз застосування МГЕ та МПГЕ при знаходженні напружено-деформованого стану у кусково-однорідних тілах
свідчить про можливість поєднання цих методів, зокрема, на краю області варто використовувати приграничні елементи, а на
лінії контакту -- можна граничні. При цьому слід пам'ятати, що тоді необхідно створювати одне програмне забезпечення
(ПЗ) для інтегрування по внутрішніх і приграничних елементах, а інше -- по граничних, а також те, що при обчисленні
поверхневих зусиль МГЕ з'являються інтеграли в сенсі Коші, а це вимагає попереднього аналітичного виділення особливостей
(головного значення). Тому дослідникам, які прагнуть уніфікувати своє ПЗ та бажають обмежитись тільки числовим
інтегрування краще користуватись тільки МПГЕ.
